\makeabstract
    {
    Synthesis of Poly(Carbonate-Ethers) from Organocatalytic, Ring-Opening Polymerization
    }
    {
    Jacob Dinzeo\textsuperscript{1}, Trinity Lee\textsuperscript{1}, Christopher B. Durr\textsuperscript{1}, Emily Ness\textsuperscript{1}
    }
    {
    \textsuperscript{1} Department of Chemistry, Amherst College, Amherst, MA
    }
    {
    In the medical field, poly(ethylene Glycol) (PEG) has been widely used as a water-soluble, biocompatible polymer material. However, PEG is nondegradable, leading to bioaccumulation and health hazards. Therefore, researchers have turned to degradable alternatives such as poly(carbonate-ethers) synthesized by organocatalytic, ring-opening polymerization. This work expands upon poly(carbonate-ether) synthesis by investigating the coupling of glycerol carbonate (GC) with small molecules and drugs prior to polymerization, yielding degradable, water-soluble polymers that enable easy and diverse functionalization via esterification. Additionally, we explored the synthesis of poly(carbonate-ethers) into bottlebrush polymers. The high density of functional side chains within bottlebrush architectures provides opportunities for increased drug loading and tunable viscosity, potentially enabling more versatile and effective drug delivery systems.
    }
    {
    GC was coupled with the desired small molecule or drug at room temperature for 24 hours using EDC-HCL and DMAP as coupling reagents. The resulting monomer was subjected to ring-opening copolymerization with ethylene carbonate (EC) (20\% GCX: 80\% EC) at 160 {\textcelsius} for 4 hours using TBD as an organocatalyst and benzyl alcohol as an initiator. For the bottlebrush polymers, norbornene-initiated macromonomers were subjected to ring-opening metathesis polymerization (ROMP) using a Grubbs 3rd-generation catalyst overnight. All structures were confirmed by 1H NMR and, when suitable crystals formed, by SCXRD. All reactions were handled and synthesized in an anhydrous, O2-free glovebox.
    }
    {
    We successfully coupled GC with BA to form GCBA, and with naproxen (N) to form GCN, as evidenced by 1H NMR and, when relevant, SCXRD. In addition, 1H NMR supports the successful copolymerization of GCBA and GCN with EC to form poly(EC-co-EO-co-GCBA-co-GOBA) and poly(EC-co-EO-co-GCN-co-GON), respectively. Finally, a norbornene-initiated GCBA macromonomer was successfully synthesized, as shown by 1H NMR. However, no bottlebrush polymers were successfully synthesized. 
    }
    {
    We examined the functionalization of GC with small molecules and drugs prior to copolymerization with EC. Additionally, we explored the synthesis of norbornene-initiated poly(carbonate-ether) bottlebrush polymers. Together, functionalized poly(carbonate-ethers) and their bottlebrush derivatives provide a promising platform for developing biodegradable, water-soluble, and highly functionalizable alternatives to PEG for biomedical applications.
    }

    \newpage
    